\section{The compact branching law from polynomial representations}
\label{cb-section}

We prove the branching law needed both for the integral lowering map
and for the terminal radial operator.  The proof uses the polynomial
modules and their duality from Section~\ref{coef:section}.  In
particular, multiplicity one and the precise interlacing inequalities
will be proved, not assumed as additional representation-theoretic
inputs.

For a partition $\nu=(\nu_1,\ldots,\nu_r)$, use the polynomial
module $V_\nu=\nabla_\nu(\mathbb C)$ on an $r$ by $r$ matrix.
Zero last columns may be omitted: column degree zero means that
the polynomial does not involve that column.  We use the representation
\[
             (\rho(g)F)(X)=F(g^tX),\qquad g\in\operatorname{GL}_r(\mathbb C).
\]
The transpose makes $\rho(g)\rho(h)=\rho(gh)$.
Its infinitesimal operators are
\[
 E_{ab}=\sum_s x_{as}\frac{\partial}{\partial x_{bs}},
 \qquad
 [E_{ab},E_{cd}]=\delta_{bc}E_{ad}-\delta_{da}E_{cb}.
 \tag{CB1}
\]
For the Hermitian Fock form, linear in its first argument,
$E_{ab}^*=E_{ba}$.  This follows directly on monomials,
or by Lemma~\ref{coef:adjoint}.  The restriction of $\rho$ to
$U(r)$ is unitary.  The highest polynomial $H_\nu$ from
Section~\ref{coef:section} has row weight $\nu$ and satisfies
$E_{ab}H_\nu=0$ for $a<b$.

For a dominant integral tuple $\lambda_1\ge\cdots\ge\lambda_r$,
put $s=\lambda_r$ and define
\[
 V_\lambda=\det^{\,s}\otimes V_{\lambda-s(1,\ldots,1)}.
 \tag{CB2}
\]
Negative $s$ causes no difficulty on $\operatorname{GL}_r$ or $U(r)$.
The determinant character has absolute value one on the latter.
Its infinitesimal action adds $s$ to every diagonal weight.

\subsection{Highest weights and irreducibility}

\begin{lemma}\label{cb-polynomial-irreducible}
The representation $V_\lambda$ is irreducible, and its highest
weight is $\lambda$.
\end{lemma}
\begin{proof}
First suppose that $\lambda=\nu$ is a partition.  Theorem
\ref{coef:integral-duality} shows that the divided row operators
applied to $H_\nu$ span $\nabla_\nu(\mathbb Q)$.  The same is
true over $\mathbb C$.  Indeed these spaces have a fixed finite
monomial support, and their equations have rational coefficients:
column homogeneity and $E^{\rm col}_{ij}F=0$ for $i<j$ give their
defining equations.  In characteristic zero the latter differential
equation is equivalent to invariance under the corresponding column
addition, by the finite Taylor expansion.  Taking a kernel of a
finite rational matrix commutes with extension to $\mathbb C$.

Every divided off-diagonal row operator satisfies
$D_{ab}^{(k)}=E_{ba}^k/k!$;
the divided diagonal operator is $\binom{E_{aa}}k$ on each row
degree.  Thus $H_\nu$ is cyclic for the operators (CB1).
Suppose $W$ is a $U(r)$-invariant subspace of $V_\nu$ and let
$P$ be its orthogonal projection.  Its orthogonal complement is
invariant too, so $P$ commutes with $U(r)$ and hence with its
complexified infinitesimal operators (CB1).  It preserves row
weights.  Lemma~\ref{coef:extremal-weight} says that the row
weight $\nu$ space is exactly $\mathbb C H_\nu$.
Consequently $PH_\nu=cH_\nu$, and $P^2=P$ gives $c=0$ or $1$.
Cyclicity then gives $P=0$ or $1$.  This proves irreducibility.
Twisting by the one-dimensional determinant character preserves
irreducibility and shifts the highest weight as in (CB2).
\end{proof}

\begin{lemma}\label{cb-highest-classification}
Every irreducible finite-dimensional continuous unitary
representation of $U(r)$ is isomorphic to exactly one $V_\lambda$,
where $\lambda$ is a dominant integral tuple.  A highest vector
of weight $\lambda$ determines an intertwiner from $V_\lambda$
uniquely up to its prescribed nonzero scalar value on $H_\lambda$.
\end{lemma}
\begin{proof}
Smoothness follows from the convolution argument of
Lemma~\ref{rp-unitary-derivative}; that argument uses no
classification theorem.  The diagonal torus acts by commuting
unitary operators, which have a common orthonormal eigenbasis.
Each continuous character of a circle is $e^{it}\mapsto e^{imt}$
for an integer $m$: differentiating its smooth character gives
$e^{ict}$ for a real $c$, and period $2\pi$ gives $c\in\mathbb Z$.
Thus the torus weights are integral tuples, and only finitely
many occur.

Choose a nonzero vector $v$ of lexicographically maximal weight
$\lambda$.  For $a<b$, the operator $E_{ab}$ increases that
weight by $e_a-e_b$; the commutator with the diagonal operators
proves this assertion.  Hence $E_{ab}v=0$.  Moreover,
\[
 \|E_{ba}v\|^2
   =\langle E_{ab}E_{ba}v,v\rangle
   =(\lambda_a-\lambda_b)\|v\|^2\ge0.
\]
Here we used $E_{ab}^*=E_{ba}$ and
$[E_{ab},E_{ba}]=E_{aa}-E_{bb}$.
It follows that $\lambda$ is dominant.

We give the required uniqueness argument, including its
existence assertion.  In any unitary representation, the
Gram matrix of the vectors $Pv$, where $P$ ranges over words
in the $E_{ab}$, is determined by $\lambda$ and $\|v\|^2$.
To see this, write
\[
                    \langle Pv,Qv\rangle
                         =\langle Q^*Pv,v\rangle.
\]
Use (CB1) to order every word with lowering operators
$E_{ab}$, $a>b$, on the left, diagonal operators in the
middle, and raising operators, $a<b$, on the right.
This process terminates: swapping an adjacent pair removes
an inversion, and the commutator term has smaller word length.
Induct first on length, and then on the number of inversions.
A nonempty raising part annihilates $v$.  The diagonal part
acts by its polynomial in $\lambda$.  A nonempty lowering
part changes the weight by a nonzero sum of negative positive
roots and is therefore orthogonal to $v$.  Only the term
with neither raising nor lowering operators contributes.
This computes the Gram matrix using just $\lambda$.

Apply this calculation to $v$ and to the highest vector in
$V_\lambda$, normalizing both to have norm one.
The rule $PH_\lambda\mapsto Pv$ is well-defined: a linear
relation has squared norm zero on one side exactly when it
has squared norm zero on the other.  It is an isometry and
intertwines all $E_{ab}$.  Its image is the cyclic subspace
generated by $v$.  The group $U(r)$ is connected and is
generated by exponentials of its Lie algebra; explicitly,
unitary diagonalization writes any unitary matrix as an
exponential of a skew-Hermitian matrix.  The map is therefore
a group intertwiner.  In an irreducible representation the
cyclic subspace is the whole space.  Its uniqueness follows
from cyclicity of $H_\lambda$.

Finally the same normal ordering shows that every weight in
$V_\lambda$ is obtained from $\lambda$ by subtracting a sum
of positive roots: the module is cyclic and every word on
its highest vector reduces to lowering words.  A nonempty
such sum strictly decreases the lexicographic weight.
Thus $\lambda$ is its unique maximal weight.  An isomorphism
preserves weights, so different $\lambda$ give nonisomorphic
representations.
\end{proof}

\subsection{The highest vectors on restriction}

Let $N=n+1$ with $n\ge1$, and first let
$\lambda=(\lambda_1,\ldots,\lambda_n,0)$ be a partition.
Omit its unused last column and write the variable matrix as
$X=\left(\begin{smallmatrix}A\\z\end{smallmatrix}\right)$,
where $A$ is $n$ by $n$ and $z=(z_1,\ldots,z_n)$.
Let $\Delta_i$ be the first $i$ by $i$ minor of $A$, and let
$S_i$ be the minor with rows $1,\ldots,i-1,N$ and columns
$1,\ldots,i$.  Thus $S_1=z_1$.

\begin{lemma}\label{cb-joint-highest}
The vectors in $V_\lambda$ which have top-row weight
$\mu=(\mu_1,\ldots,\mu_n)$ and are annihilated by
$E_{ab}$ for $1\le a<b\le n$ form a space of dimension
one if
\[
                  \lambda_i\ge\mu_i\ge\lambda_{i+1}
                                      \qquad(1\le i\le n),
 \tag{CB3}
\]
and a zero space otherwise.  In the first case a generator is
\[
 \xi_{\lambda,\mu}(X)
       =\prod_{i=1}^n\Delta_i(X)^{\mu_i-\lambda_{i+1}}
                           S_i(X)^{\lambda_i-\mu_i}.
 \tag{CB4}
\]
Its last-row degree is $|\lambda|-|\mu|$.
\end{lemma}
\begin{proof}
The assertion that $E_{ab}F=0$ is equivalent, by finite
Taylor expansion, to invariance under adding a multiple
of top row $a$ to top row $b$.  These additions generate
the lower unipotent row subgroup.  By definition, $F$
is also invariant under the upper unipotent column subgroup.

On the nonempty open set $\prod_i\Delta_i\ne0$, Gaussian
elimination gives $A=LDU$, with $L$ lower unipotent,
$U$ upper unipotent, and $D=\operatorname{diag}(d_1,\ldots,d_n)$.
First multiply the columns by $U^{-1}$ and then the top
rows by $L^{-1}$.  Invariance gives
\[
                         F(A,z)=F(D,zU^{-1}).
 \tag{CB5}
\]
The row and column degrees force
\[
 F(D,w)=c\prod_i d_i^{\mu_i}w_i^{\lambda_i-\mu_i}.
 \tag{CB6}
\]
If one of these exponents is negative, the restriction is
zero, and (CB5) forces $F=0$.  In particular a nonzero
vector has $\lambda_i-\mu_i\ge0$.

On this diagonal matrix one has
$\Delta_i=d_1\cdots d_i$ and
$S_i=d_1\cdots d_{i-1}w_i$.  The Laurent product in
(CB4), allowing temporarily negative exponents
$\mu_i-\lambda_{i+1}$, restricts to the monomial in
(CB6).  It is itself invariant under both unipotent
groups.  Therefore (CB5) gives $F=c\xi_{\lambda,\mu}$
as a rational function.

The polynomials $\Delta_i,S_i$ are pairwise nonassociate
irreducibles.  Here is the elementary determinant fact
needed for this assertion.  A generic determinant is
irreducible by induction on its size: as a polynomial
in its last row it is linear, and its coefficients are
the maximal cofactors of the other rows.  Those cofactors
are irreducible by induction and no two are associates,
since they involve different column variables.  Their
greatest common divisor is one.  The polynomial is
therefore primitive and is irreducible over the fraction
field of the other-row variables, where it is a nonzero
linear polynomial.  Gauss's lemma proves irreducibility.
Size one is a variable.  Adjoining additional independent
variables preserves this conclusion.  Finally different
sizes have different degrees, and $S_i$ involves the last
row whereas $\Delta_i$ does not; hence none of the
displayed minors are associates.

Unique factorization now implies
$\mu_i-\lambda_{i+1}\ge0$ whenever $F\ne0$.
This proves the necessity of (CB3) and the upper bound
one on the dimension.  If (CB3) holds, (CB4) is a
nonzero polynomial.  Each factor is invariant under
the indicated column additions and top-row additions.
Its column degrees are $\lambda_i$, and its top-row
degrees telescope to $\mu_i$.
Each $S_i$ has last-row degree one, so the remaining
degree is $\sum_i(\lambda_i-\mu_i)$.
Thus (CB4) belongs to the required space and proves
the assertion.
\end{proof}

\begin{theorem}[Compact branching law]\label{ib-compact-branching}
For $N\ge2$ and a dominant integral highest weight
$\lambda=(\lambda_1,\ldots,\lambda_N)$, restriction
of $V_\lambda$ to $U(N-1)\times U(1)$ is the orthogonal
direct sum
\[
 V_\lambda\big|_{U(N-1)\times U(1)}
     \simeq\bigoplus_{\lambda_i\ge\mu_i\ge\lambda_{i+1}}
                           V_\mu\otimes t^{\,|\lambda|-|\mu|}.
 \tag{CB7}
\]
Each summand occurs once.  The character of the last
$U(1)$ is displayed, not suppressed.  For $N=1$ the
restriction to the trivial group has one one-dimensional
summand, of last-circle weight $\lambda_1$.
\end{theorem}
\begin{proof}
First subtract $\lambda_N$ from every weight, so that
the polynomial model of Lemma~\ref{cb-joint-highest}
applies.  A finite-dimensional unitary representation
is an orthogonal sum of irreducibles: choose a minimal
nonzero invariant subspace, take its invariant orthogonal
complement, and induct on dimension.
Lemma~\ref{cb-highest-classification} identifies every
summand of the restriction to $U(N-1)$ by a highest
vector.  Lemma~\ref{cb-joint-highest} allows exactly
the interlacing weights and at most one copy of each.
Conversely (CB4) provides a highest vector for every
interlacing weight.  The Gram-matrix argument in
Lemma~\ref{cb-highest-classification} identifies its
cyclic subspace with $V_\mu$, so every listed summand
does occur.  The last-row degree in that lemma proves
the $U(1)$ character.  Inequivalent unitary summands
are orthogonal: their cross-pairing defines an
intertwiner, which is zero by irreducibility and the
distinct highest weights.

Twisting back by $\det^{\lambda_N}$ adds $\lambda_N$
to all the top weights and to the last-circle weight.
It preserves (CB3) and changes the last exponent to
$|\lambda|-|\mu|$, as stated.  For $N=1$ (CB2) is
just the character $t\mapsto t^{\lambda_1}$.
\end{proof}

This proof also specifies the multiplicity-one normalization:
in the polynomial case there is a unique intertwiner
$J_{\lambda,\mu}:V_\mu\to V_\lambda$ with
$J_{\lambda,\mu}H_\mu=\xi_{\lambda,\mu}$.
The next section will address its integral denominator;
multiplicity one alone does not assert that this
normalization preserves the dual-Weyl lattice.


